\section{Ablations}

The knockouts settle it. The budget is better spread. Ablations run in three settings
spanning a CNN policy, a state diffusion policy and an image diffusion policy, in
the order GridWorld (image), Push-T (state), Door (image). Lift is excluded for the
reason given above: at the ceiling a null result is unreadable. We report the three
values and their signs. The supplementary gives the full programme with margin
retained per setting.

\paragraph{Removing the partition costs success without costing information.}
Deleting the mode structure but keeping the loss signal, so each round greedily
corrects the highest-loss failure, is the most damaging knockout, costing
$-2.2 / -4.1 / -6.8$ points, a mean of $4.37$, and a mean of $-53.2$ per cent of the
margin retained, beneath its own best baseline on Push-T (92.0 against 94.1) and Door
(81.8 against 84.2). Per-demonstration information gain does not fall with it
($+0.02 / +0.16 / -0.13$, mean $+0.02$). Greedy worst-loss selection acquires
demonstrations that are individually informative and jointly redundant, so the
partition buys coverage of the failure distribution and not a better
demonstration. The knockout removes the whole allocation stack, since the
descriptor and the memory have nothing to operate on once modes are gone.

\paragraph{The advantage lies in the allocation.}
Uniform-random allocation over the recorded failures (A2) reaches $89.1 / 82.3 / 80.0$
against DISEIL's $91.3 / 96.1 / 88.6$, below the strongest gated baseline on both robot
settings, and promoting the deterministic fallback to the whole method costs $3.27$
points, so neither replaying failures nor the heuristic accounts for the result. That
fallback knockout keeps the partition and reverts only the placement, the nearest
isolation of \emph{where} available here, but it also removes the prescription
model, so the two are not separated. Removing the constraint graph costs $2.37$ points and
drives the fallback rate to 27 to 35 per cent of rounds, and bridging placement costs
$1.27$; the prescription model, the vision-language model and the cluster memory cost
$1.33$, $1.33$ and $0.73$, gaps comparable to the seed standard error of the full
run. Replacing the prescription model with a heuristic leaves the Push-T margin at
$0.1$ points, inside that standard error, so a model-free deployment retains the
result on two of the three settings only; dropping the vision-language model as well
was not measured. Across budgets the margin is $+9.07$ at $B{=}10$, $+2.87$ at
$B{=}20$ and $+2.83$ at $B{=}40$, narrowing because the baselines catch up from a
lower start, though DISEIL at $B{=}10$ does not reach what the strongest baseline
attains at $B{=}20$.

\section{Limitations and Conclusion}

DISEIL reasons about one round and holds no representation of the training set. Its
only memory records where corrections were placed, not what the data cover, so it
cannot tell a failure the set never taught from one it already teaches six times and
is merely under-fitted. The descriptor is designed by hand and recovers cause only
where configuration determines cause, with root-cause purity per cluster at 0.84 to
0.91 scored against the prescription model's own labels, so it records agreement
inside one system. It is also read from privileged simulator state, so the image
results show the allocation transferring to an image policy, not the partition built
from pixels; outside simulation it would need a pose estimator whose error would
enter the partition. The reasoning stack costs 63 to 1{,}232 seconds per round more
than a query gate, at up to 82{,}116 tokens, though it runs only at selection time
and never in the control loop. No knockout varies \emph{where} with \emph{which}
held fixed, so the split of the margin between the two is not measured here. Every
expert here is a simulator oracle, which is what lets a budget be counted in
demonstrations at all. Outside simulation the budget is a person's time and
demonstrators are not uniformly correct \cite{belkhale2023dataquality,
brown2019trex}, where the information-gain argument fails, since a high pre-retrain
loss identifies novel data only if invalid demonstrations are excluded by
construction.

The which and the where are free only while the budget is not. DISEIL takes both
by partitioning a round's failures into modes over a geometric descriptor and
prescribing a feasibility-checked configuration, and it holds the best or
joint-best mean success rate in every setting reported here. Removing the
allocation stack leaves per-demonstration information gain untouched while success
falls by a mean of 4.37 points across the three settings ablated: the
demonstrations stayed informative and stopped covering the failure distribution.
The components of the where each cost between 0.7 and 2.4 points, inside the seed
standard error of the full run, so this study does not measure what the
prescription is worth on its own. Choice of gate moves the result by as much as 19
points, on Wipe (image), yet no gate exceeds DISEIL in any of the ten settings.
What a round buys turns on which failure it corrects, not only on which scalar
opens the gate.
